% BHU PYQ -- Professional Exam Template
\documentclass[11pt,a4paper]{article}

\usepackage[a4paper,top=2.5cm,bottom=2.5cm,left=2.8cm,right=2.8cm,headheight=14pt]{geometry}
\usepackage{amsmath,amssymb}
\usepackage[T1]{fontenc}
\usepackage[utf8]{inputenc}
\usepackage{lmodern}
\usepackage{microtype}
\usepackage{fancyhdr}
\usepackage{enumitem}
\usepackage{hyperref}
\usepackage{lastpage}
\usepackage{parskip}

\hypersetup{colorlinks=true,urlcolor=black,linkcolor=black,
  pdftitle={B.Sc. (Hons.) Zoology Sem-V ZOB-504 -- BHU PYQ}}

\pagestyle{fancy}
\fancyhf{}
\renewcommand{\headrulewidth}{0.4pt}
\renewcommand{\footrulewidth}{0.4pt}
\fancyhead[L]{\small   \textit{Banaras Hindu University}}
\fancyhead[R]{\small   \textit{Zoology Paper: ZOB-504}}
\fancyfoot[C]{\small Page  \thepage\ of \pageref{LastPage}}


\newlist{parts}{enumerate}{3}
\setlist[parts,1]{label=(\alph*),leftmargin=*,itemsep=4pt,topsep=3pt}
\setlist[parts,2]{label=(\roman*),leftmargin=*,itemsep=2pt,topsep=2pt}
\setlist[parts,3]{label=(\arabic*),leftmargin=*,itemsep=2pt,topsep=2pt}

\newcommand{\pts}[1]{\hfill   \textit{[#1]}}


\begin{document}


\begin{center}
  {\large\bfseries BANARAS HINDU UNIVERSITY}\\[2pt]
  {
\normalsize B.Sc. (Hons.) Semester V Examination, 2016--17}\\[6pt]
  \rule{0.8\linewidth}{0.5pt}\\[6pt]
  {\large\bfseries Zoology}\\[4pt]
  {\large\bfseries Paper No. ZOB-504 : Biotechniques}\\[4pt]
  {
\normalsize Time: 3 hours \hfill Full Marks: 70}
\end{center}

\rule{\linewidth}{0.4pt}

\smallskip



\noindent   \textit{Note: Write your Roll No. at the top immediately on the receipt of this question paper. The figures in the right-hand margin indicate marks. Answer six questions including Question 1 which is compulsory.}

\bigskip



\noindent   \textbf{Question 1.} Answer the following parts:

\begin{parts}
  \item Select and write the best options in the following:\pts{$4  \times \frac{1}{2} = 2$}
  
\begin{parts}
    \item In SDS-PAGE, acrylamide determines
    
\begin{parts}
      \item pH of Gel
      \item pore size
      \item electrophoretic mobility
      \item ionic strength
    \end{parts}

    \item Resolving power of microscope is dependent upon
    
\begin{parts}
      \item numerical aperture
      \item eye piece
      \item magnification
      \item body tube
    \end{parts}

    \item Swelling followed by rupture of cells   \textit{in vitro} is due to change in
    
\begin{parts}
      \item pH
      \item osmolarity
      \item temperature
      \item ionic strength
    \end{parts}

    \item Substances get separated during centrifugation due to
    
\begin{parts}
      \item centrifugal force
      \item centripetal force
      \item gravitational force
      \item relative centrifugal force
    \end{parts}
  \end{parts}

  \item State whether the following statements are True or False. Justify your answer in brief:\pts{$6  \times 2 = 12$}
  
\begin{parts}
    \item DMSO is not required in culture medium during cryopreservation.
    \item The electrochemical potential between two liquids is measured by pH meter.
    \item Geiger-Muller counter can detect only beta particles.
    \item Serum is not an important requirement for adherent cell culture.
    \item SEM produces more resolution than TEM.
    \item Ethidium bromide is commonly used to detect proteins in PAGE.
  \end{parts}

  \item Define the following:\pts{$4  \times 1 = 4$}
  
\begin{parts}
    \item Electrophoresis
    \item Antibody
    \item Limit of resolution
    \item Cell fractionation.
  \end{parts}

  \item Expand the following:\pts{$4  \times \frac{1}{2} = 2$}
  
\begin{parts}
    \item TEMED
    \item FISH
    \item EIA
    \item IHC.
  \end{parts}
\end{parts}

\pagebreak

\medskip


\noindent   \textbf{Question 2.} Describe in detail with labeled diagram the principle and application of phase-contrast microscope.\pts{10}

\medskip


\noindent   \textbf{Question 3.} Describe the general composition of cell culture media. Comment upon their uses and significance.\pts{10}

\medskip


\noindent   \textbf{Question 4.} What is radioactivity? Describe the measurement and safety aspects of radioactivity.\pts{10}

\medskip


\noindent   \textbf{Question 5.} What is microtomy? Enumerate the procedure for the preparation of histological HE stained slides.\pts{10}

\medskip


\noindent   \textbf{Question 6.} Write down the principle of centrifugation. Describe the various types of centrifuges used for the separation of biomolecules.\pts{10}

\medskip


\noindent   \textbf{Question 7.} Write down the procedure of DNA extraction using mammalian tissue. Describe the principle and application of DNA-gel electrophoresis.\pts{10}

\medskip


\noindent   \textbf{Question 8.} Write short notes on the following:\pts{$2  \times 5 = 10$}

\begin{parts}
  \item Spectrophotometer
  \item Cryopreservation.
\end{parts}

\vfill
\rule{\linewidth}{0.4pt}

\begin{center}\small
\href{https://bhu-exam.vercel.app/}{Click Here}\\[4pt]\href{https://me-aryan.vercel.app/}{Click Here}\\[4pt]\href{https://quashan-me.vercel.app/}{Click Here}
\end{center}

\end{document}
